\chapter{从理论对象到可训练模型：表示、架构与优化}
\label{th:representation}

\section{概率建模与神经网络结构分开选择}

DDPM、分数或流匹配规定目标与过程，不规定只能使用 U-Net。
同一个 DiT 可输出噪声、旋转参数或速度，但权重、调制与损失要按目标训练。
反过来，相同的 $u_\theta$ 记号也可能由卷积、Transformer 或图网络实现。
架构影响归纳偏置与计算，不能由网络名称反推出概率路径。

U-Net 通过下采样、瓶颈、上采样与跨层跳接结合局部纹理和全局语义。
时间 embedding 常进入残差块的尺度与偏置；
文本可经 cross-attention 与中间特征交互。
跨层跳接传递同一样本的特征，不是把真实干净目标泄漏给生成器。

DiT 将空间或时空位置投影为 token，经过注意力与前馈网络。
给定 $Q=XW_Q,K=XW_K,V=XW_V$，
\[
Y=\operatorname{softmax}(QK^\top/\sqrt{d_h}+A)V.
\]
时间/条件可通过自适应归一化调制，例如
$\operatorname{AdaLN}(h;c)=(1+\gamma(c))\operatorname{LN}(h)+\beta(c)$，
再用门控缩放残差。具体零初始化策略需要防止梯度路径被同时切断。

文本与图像分别投影后联合自注意力，与仅由图像 query 读取文本的交叉注意力不同。
可见性矩阵 $A$ 决定哪些位置交换信息；
padding、条件缺失、因果控制和局部编辑不应合并为一个含糊的 mask。

\section{latent diffusion 的两阶段目标}

令 $z=E_\phi(x)$，扩散模型学习编码后的分布。
解码结果的质量受到 codec 重建能力限制：压缩丢掉的小字与细纹，
即便潜变量生成完全准确也不能自动找回原始细节。
压缩可以显著减少 token 与计算，也会改变目标的统计尺度与可表示性。

概率化模型若有 $p_\psi(x\mid z)$ 与 $p_\theta(z)$，则
\[
\log p(x)\geq\E_{q_\phi(z\mid x)}\log p_\psi(x\mid z)
 -\KL(q_\phi(z\mid x)\Vert p_\theta(z)).
\]
而实际常用的确定编码、感知重建与潜空间去噪，未必直接优化这条像素 ELBO。
潜空间精确密度不能不加 codec/量化模型就当作原像素的精确似然。

需记录编码后缩放、平移、通道标准差、采样后验还是使用均值。
训练用 $z/s$、推理忘记乘回 $s$，可能使模型在完全错误的 SNR 区间生成。
不同 codec 即使 shape 相同也不意味着通道语义和权重兼容。

\section{损失平均、缺失模态与梯度路径}

对样本 $b$，有效元素权重为 $W_{bi}\geq0$，推荐先定义清楚
\[
\ell_b=\frac{\sum_iW_{bi}\|\hat y_{bi}-y_{bi}\|^2}{\sum_iW_{bi}},
\]
再对有效样本平均。这里 $i$ 若已包含标量通道，不能再重复除以通道数。
若用 token mask，需要先说明广播后的分母。
样本长度不同，按全部 token 一起平均与按样本平均优化不同的数据权重。
没有有效监督的样本应排除，而非加小数后贡献一个零损失。

冻结参数只表示不更新该权重。
若前面的新模块需要经过后续冻结层获得梯度，仍需保留这段计算图。
把整个冻结主干包在 no-grad 中会截断新模块训练。
LoRA 降低可训练状态量，并不消除长序列前向与反传激活成本。

\section{预条件、精度与噪声课程}

网络输入输出的尺度会影响优化。
EDM 预条件是一种明确的缩放设计，时间分布与 SNR 权重则决定不同难度任务的资源。
大噪声多时偏重全局布局，小噪声多时偏重细节回归；
这个描述不意味着两类能力可完全独立训练。

BF16/FP16 适合矩阵计算，损失累计、方差、极端指数和某些位置频率
可能需要更高精度。梯度裁剪与 EMA 是稳定/平均策略，
不是对错误目标或错误 scheduler 的补救证明。
改变分辨率、长宽比、时长或条件数，会改变信号尺度与 token 分布，
仅因线性层参数 shape 不变不能断言训练分布已适配。

\section{理论完备不能代替数据与可复现性}

真实训练要保存模型、优化器、学习率调度、随机状态、数据游标、
噪声采样规则、codec 版本与参数化。
否则“继续训练”可能只是加载权重后开启一个不同的优化过程。
验证需要固定数据来源与去重规则，避免同主体或同原片的变体穿过数据划分。
训练损失、生成质量与条件遵循要分别记录。

\begin{theorytakeaway}
表示与网络决定可学习范围，训练协议决定实际优化分布。
概率理论说明应学什么，不能省略 codec、缩放、梯度和数据版本这些实现前提。
\end{theorytakeaway}
\threading{主参考的理论框架\cite{theorybook}与 EDM 的尺度设计\cite{edm}。
本章为通用架构接口说明，不推断特定闭源模型，也不改动后面的 Seedance 专题。}
